\title{KronA: Parameter Efficient Tuning with Kronecker Adapter}

\begin{abstract}
Fine-tuning a Pre-trained Language Model (PLM) on a specific downstream task has been a well-known paradigm in Natural Language Processing. However, with the ever-growing size of PLMs, training the entire model on several downstream tasks becomes very expensive and resource-hungry. Recently, different Parameter Efficient Tuning (PET) techniques are proposed to improve the efficiency of fine-tuning PLMs. One popular category of PET methods is the low-rank adaptation methods which insert learnable truncated SVD modules into the original model either sequentially or in parallel. However, low-rank decomposition suffers from limited representation power. In this work, we address this problem using the Kronecker product instead of the low-rank representation. We introduce KronA, a Kronecker product-based adapter module for efficient fine-tuning of Transformer-based PLMs. We apply the proposed methods for fine-tuning T5 on the GLUE benchmark to show that incorporating the Kronecker-based modules can outperform state-of-the-art PET methods.
\end{abstract}

\section{Introduction}
Large PLMs are used as a backbone model in a variety of NLP tasks to achieve state-of-the-art results \cite{devlin-etal-2019-bert,radford2019language}.  These large pre-trained models are adapted to the downstream applications either via in-context learning or fine-tuning of the model parameters. In-context learning imposes substantial memory and computational overhead during inference as all the training examples have to be processed for each sample \cite{liu2022few}. On the other hand, full Fine-Tuning (FT) the entire model provides both less inference latency and improved accuracy. However, as these models become larger, full fine-tuning of their parameters becomes more challenging. Additionally, one has to store an entire model checkpoint for each downstream application, which makes deployment and switching between different tasks extremely inefficient. 

To address these challenges, several works have proposed to insert a small number of trainable parameters while freezing most (or even all) of the pre-trained model parameters. This significantly reduces the memory and computation requirements for fine-tuning. Furthermore, instead of storing one copy of the entire model, a small set of tuned parameters can be stored for each task. We refer to these methods as PET methods. 

Among the PET methods, soft prompts \cite{li-liang-2021-prefix,lester2021power} prepend trainable parameters to the input of the layers. The increase in length of the embedding layers leads to a significant computation overhead during the inference. 

In another category of the PET methods, adapter modules are inserted  \cite{pmlr-v97-houlsby19a,karimi2021compacter,he2022towards} into the Transformer. Adapters are low-rank modules that are composed of an up projection followed by a down projection. One limitation of these approaches is that they increase the computational overhead and the latency during the inference which makes them inefficient for latency-critical scenarios. 

Therefore, Low Rank Adaption (LoRA) \cite{hu2021lora} was developed, which also uses extra low-rank modules as the trainable parameters. However, once fine-tuned, the task-specific parameters can be merged with the original pre-trained model weights,  making the latency and energy requirements for inference, intact. Despite of fast inference, usually LoRA suffers from an accuracy drop compared to full fine-tuning. This is because of the strong assumption imposed by its low-rank structure for task-specific updates. 

Kronecker product decomposition is another factorization method that does not rely on the low-rank assumption. This powerful decomposition method, when used for model compression, has proven to outperform low-rank factorization methods \cite{thakker2019pushing,hameed2021convolutional}. It has also been successfully used for the compression of Transformer-based language models \cite{tahaei2021kroneckerbert,edalati2021kronecker}.

Inspired by the ubiquitous success of Kronecker decomposition, in this work we replace low-rank decomposition in LoRA with Kronecker product decomposition to develop the Kronecker Adapter (KronA). We show that this simple modification can improve the accuracy without increasing the inference latency. Also for applications where an increase in latency is tolerable, we propose to use $\text{KronA}^\text{B}$. This module is a version of KronA developed to be utilized in parallel to Feed-Forward Network (FFN) modules and achieves significant improvements over full fine-tuning on the General Language Understanding Evaluation (GLUE) benchmark \cite{wang-etal-2018-glue}. In addition, when a proposed learnable residual connection is added to $\text{KronA}^\text{B}$, $\text{KronA}^\text{B}_\text{res}$ is developed to achieve even better results.

We evaluated our methods on the GLUE benchmark \cite{wang-etal-2018-glue} to study the impact of Kronecker product on the performance. To summarize, our contributions are:

\begin{itemize}
    \item Proposing KronA, a \underline{Kron}ecker \underline{A}dapter module that can be inserted in parallel to the weight matrices and is suitable for latency-critical scenarios.
    \item Using the KronA module in parallel to the FFN module ($\text{KronA}^\text{B}$) along with a learnable residual connection ($\text{KronA}^\text{B}_\text{res}$) to further improve the accuracy at the cost of increased inference latency.
    \item Providing evaluation of the methods in comparison to the state-of-the-art in terms of the GLUE score, training time and inference latency.
\end{itemize}

\section{Related Works and Baselines} \citep{zaken2021bitfit} proposed freezing the weights and tuning only biases or a subset of biases in the PLM to fine-tune it on the downstream tasks. This technique, called BitFit, is parameter-efficient and fast, but it usually cannot achieve good performance compared to the state-of-the-art methods.

\cite{pmlr-v97-houlsby19a} introduced Adapters as a PET method. In this method, the entire parameters of a PLM are frozen and some trainable modules, called Adapters, are inserted after the FFN or attention blocks, sequentially. Every Adapter module has a down projection, a non-linear function and an up projection in addition to a residual connection. \cite{he2022towards} developed Parallel-Adapter (PA) which outperforms Adapter. PA is inserted in parallel to the original PLM module and has a scaling factor, but the residual connection inside the module is removed since the PLM module has a residual connection (Figure \ref{fig:1}.c). \cite{he2022towards} also introduced a unified view over PET methods and combined some of the techniques like Prefix tuning \cite{li-liang-2021-prefix} and PA.

In \cite{karimi2021compacter}, a modified version of the sequential adapter is used for PET, named Compacter. In Compacter, Kronecker product of multiple Kronecker factors are added to reconstruct the module's weight matrix. Each Kronecker factor itself is a result of matrix multiplication of two sub-factors. Compacter achieves good results on GLUE, but it is notably slow in both training and inference phases.

LoRA \cite{hu2021lora} inserts trainable modules which are made from a down projection and an up projection in different parts of a PLM. It is worth mentioning that the authors recommended to insert the LoRA modules in parallel to the query and value matrices.  During the training, the PLM weights are frozen and only the LoRA modules are tuned. During the inference, the LoRA weights are merged with the original weight matrices of the PLM. Therefore, in contrast to the Adapter and Compacter, LoRA does not increase the inference time. 

Since LoRA, PA and Compacter have outperformed most of the existing baselines such as Pfeiffer-Adapter \cite{pfeiffer2020adapterfusion}, AdapterDrop \cite{ruckle-etal-2021-adapterdrop}, VGLM-Adapter \cite{lin-etal-2020-exploring} Prompt tuning \citep{lester-etal-2021-power}, and Prefix tuning \cite{li-liang-2021-prefix}, we considered these works as baseline for comparison in this work.

\section{Methodology}
\label{sec:meth}
\subsection{Kronecker Product} Kronecker product is an operation on two input matrices ($\mathbf{A} \in \mathbb{R}^{a_1 \times a_2}$ and $\mathbf{B} \in \mathbb{R}^{b_1 \times b_2}$) that results a block matrix ($\mathbf{W} \in \mathbb{R}^{w_1 \times w_2}$). In $\mathbf{W}$, each block $(i, j)$ is equal to the multiplication of the element $a_{i,j}$ and the matrix $\mathbf{B}$. Also, the shape of the resulting matrix is $(w_1, w_2)$, where $w_1 = a_1 \times b_1$ and $w_2 = a_2 \times b_2$. Equation \ref{eq:kron} shows how the Kronecker product works. for more details see  \cite{Henderson1983OnTH}.

\begin{equation}
    \mathbf{W} = \mathbf{A}\otimes\mathbf{B} = \begin{bmatrix}
  a_{11} \mathbf{B} & \cdots & a_{1n}\mathbf{B} \\
             \vdots & \ddots &           \vdots \\
  a_{m1} \mathbf{B} & \cdots & a_{mn} \mathbf{B}
\end{bmatrix}
\label{eq:kron}
\end{equation}

The Kronecker product has some interesting features that makes it suitable for PET. First, it is not rank deficient. In other words, unlike the low-rank down-projection used in LoRA and Adapter, Kronecker product decomposition maintains the rank of the input matrix. Second, to speed up and reduce the required number of FLOPS, one can avoid the reconstruction of $\mathbf{W}$ and use Equation \ref{eq:kronecker_equivalance} to calculate the output of a Kronecker module. More precisely, multiplying an input vector $\mathbf{x} \in \mathbb{R}^{d_h}$, where $d_h$ is the embedding dimension, by the Kronecker product of $\mathbf{A}$ and $\mathbf{B}$ can be performed using the following equation:

\begin{equation}
    (\mathbf{A} \otimes \mathbf{B})\mathbf{x} = \gamma(\mathbf{B}\eta_{b_2\times a_2}(\mathbf{x})\mathbf{A^\top})
\label{eq:kronecker_equivalance}
\end{equation}

where  $\mathbf{A}^\top$ is $\mathbf{A}$ transpose. $\eta_{m \times n}(\mathbf{y})$ is a mathematical operation that reshapes a vector  $\mathbf{y} \in \mathbb{R}^{mn}$  into a matrix of size $m \times n$. $\gamma(\mathbf{Y})$ is another operation that converts a matrix  $\mathbf{Y} \in \mathbb{R}^{m\times n}$ into a vector by stacking its columns. 

The following two sub-sections describe our proposed Kronecker modules. See Appendix \ref{sec:ab} for the ablation experiments justifying our design choices. These include studying the effect of initialization, the use of scale factor and non-linearity, as well as the position of the Kronecker module (parallel vs sequential). 

\subsection{KronA}
Figure \ref{fig:1}.a shows the structure of a LoRA module where $\mathbf{A}$ is the down projection and $\mathbf{B}$ is the up projection. To modify this module into KronA, the normal product is replaced by the Kronecker product. Also, LoRA projections are replaced with the Kronecker factors (See Figure \ref{fig:1}.b and Appendix  \ref{sec:app_a}). Equation \ref{eq:krona} shows how the output is generated when KronA is applied. $\mathbf{A}_k$ and $\mathbf{B}_k$ are the Kronecker factors that replaced the LoRA projections. Similar to LoRA, our KronA has a fixed scale factor, $s$, which is a hyperparameter.

\begin{equation}
\label{eq:krona}
\mathbf{Y}= \mathbf{X}\mathbf{W} + s\mathbf{X}[\mathbf{A}_k\otimes\mathbf{B}_k]
\end{equation}

KronA modules are applied in parallel to the weight matrices in a PLM during the tuning phase. Once fine-tuned, the Kronecker factors are multiplied, then scaled and merged to the original PLM weight matrix (Equation \ref{eq:krona-merge}). Therefore, similar to LoRA, KronA does not increase the inference time.

\begin{equation}
\label{eq:krona-merge}
\mathbf{W}_{\text{tuned}}=\mathbf{W}+s[\mathbf{A}_k\otimes\mathbf{B}_k]
\end{equation}

For initialization of the Kronecker factors, we observed that initializing one of the factors with zero significantly improves the results compared to initializing both factors using the Normal distribution. See Appendix \ref{sec:ab} for more details.

\subsection{$\text{KronA}^\text{B}$ and $\text{KronA}^\text{B}_\text{res}$}
Inspired by the promising performance of the PA method, we also investigate KronA module when used in parallel to the FNN block and call it KronA$^\text{B}$. The B superscript in the name means that this module is applied to the PLM \underline{b}locks, opposed to KronA which is applied to the PLM weight matrices. Note that similar to the PA method, the non-linearity in the FFN block does not allow our KronA module to be merged into the pre-trained model after fine-tuning. This imposes an increase in the inference time and computations. Equation \ref{eq:kronam} shows how KronA$^\text{B}$ works in parallel to the FFN block.

\begin{equation}
\label{eq:kronam}
\mathbf{Y} = \text{FFN}(\mathbf{X}) + s\mathbf{X}[\mathbf{A}_
k\otimes \mathbf{B}_k]
\end{equation}

To further improve the representation power, we incorporate a scaled residual connection inside the $\text{KronA}^\text{B}$ module to develop $\text{KronA}^\text{B}_\text{res}$. Scale of the residual connection ($s_{res}$) is initialized with one and tuned during the fine-tuning. Equation \ref{eq:kronamr} shows how the $\text{KronA}^\text{B}_\text{res}$ works in parallel to the FFN. Also, Figure \ref{fig:1}.c and \ref{fig:1}.d show the structure of the PA module and our $\text{KronA}^\text{B}_\text{res}$, respectively.

\begin{equation}
\label{eq:kronamr}
\mathbf{Y} =  \text{FFN}(\mathbf{X}) + s\mathbf{X}[\mathbf{A}_K\otimes \mathbf{B}_K] + s_{res}\mathbf{X}
\end{equation}

\section{Results and Discussion}
\label{sec:exp}
\subsection{GLUE Results} Table \ref{tab:1} shows\footnote{Details regarding the hyperparameters and experimental setups are mentioned in Appendix \ref{sec:app_H}.} the GLUE score of our proposed methods compared to other baselines when applied to the T5 model \cite{raffel2019exploring}. As the results show, KronA and KronA$^\text{B}$ outperform LoRA and PA as their low-rank counterparts, respectively. Also, all of our proposed modules outperform other state-of-the-art baselines on average and most of the GLUE tasks. In addition, $\text{KronA}^\text{B}_\text{res}$ which benefits from an extra learnable residual connection achieves remarkably better results. 

\subsection{Inference and Training Time} Table \ref{tab:2} shows the normalized inference delay for the discussed methods. KronA, LoRA, FT, and Bitfit do not increase the inference latency since these techniques do not add any extra parameters or computations to the model during the inference stage. Although $\text{KronA}^\text{B}$ is significantly faster than Compacter and Adapter, it is slower than PA which is as we expected; calculation of the Kronecker product is generally slower than normal matrix multiplication. Also, adding the learnable residual connection increases the latency.

Normalized training time averaged over the GLUE tasks for each technique is shown in Table \ref{tab:2}. Based on these results, the significant improvement in the accuracy of the proposed KronA module and its variants is at the expense of a slight increase in the training time compared to the low-rank counterparts like LoRA, PA, and Adapter. However, the training time increase is not remarkable and KronA modules are still significantly faster than FT.

\section{Conclusion and Future Directions}
\label{sec:con}
In this work, we developed Kronecker-based adapters by replacing the low-rank projections of the state-of-the-art PET methods with the Kronecker product. In addition to a comparison of the training and inference time, we evaluated our proposed adapter for fine-tuning T5 on the GLUE benchmark to show its superiority over the state-of-the-art baselines.

\end{document}
